\begin{thesisabstract}
\abstractparagraph{This thesis investigates the influence of soil--structure interaction on the seismic performance and preliminary fragility of a G+5 reinforced-concrete mid-size general hospital building in Nepal. Four analytical support conditions are compared: fixed base and flexible-base models representing hard, medium, and soft soil. Linear modal and response-spectrum analyses were performed in ETABS, while nonlinear static pushover capacity was evaluated using SAP2000 and the reported performance coordinates were interpreted using the FEMA~356 target-displacement framework. The preliminary fragility assessment is based on nonlinear time-history results from twelve ground motions.}

\abstractparagraph{The fundamental period increased from 0.984 s for the fixed-base model to 1.147 s for the soft-soil model. Cumulative modal mass participation exceeded 90 percent in both principal directions by Mode 6 for all models. RSx ULS roof displacement increased from 87.224 mm to 162.756 mm and RSx SLS roof displacement from 83.716 mm to 156.210 mm from fixed to soft support. Maximum RSx ULS drift increased from 0.009385 to 0.013882, and maximum RSx SLS drift increased from 0.009007 to 0.013324. RSx ULS base shear increased from 21,518.409 kN to 24,528.262 kN, while RSx SLS base shear increased from 20,652.947 kN to 23,541.745 kN. The ETABS settlement results increased from 0.000 mm to 9.209 mm over the same support sequence.}

\abstractparagraph{The SAP2000 pushover results show distinct X- and Y-direction capacity histories for all four support conditions. Eight reported fixed/hard/medium/soft X- and Y-direction performance-point coordinates were compared within the FEMA~356 target-displacement framework. The nonlinear time-history analysis includes twelve ground motions. Its mean Immediate Occupancy threshold PGA values are 0.2828, 0.2688, 0.2203, and 0.1747 g for the fixed, hard, medium, and soft models, respectively, indicating lower drift-threshold intensity as support flexibility increases. Because the available nonlinear results do not specify the analysed response direction and the Barpark/Barpak-labelled record could not be independently verified, the fragility findings are interpreted as preliminary.}

\abstractkeywords{Hospital Building, Soil--Structure Interaction, Response Spectrum Analysis, Pushover Analysis, Performance-Based Seismic Assessment, Nonlinear Time History Analysis, Seismic Fragility, Nepal.}
\end{thesisabstract}
